\documentclass[aspectratio=169]{beamer}
\usepackage{amsmath}
\usepackage{comment}
\usepackage{algorithm}
\usepackage{algpseudocode}


\title[LSE PRESS \hspace{25mm} \insertframenumber/\inserttotalframenumber]{Disagreement and Forecast Errors \\ across Firms and Time}

\author[Advanced Macroeconomics]{Hatsu Kawabata \& Tatsuro Senga}

\begin{document}
\maketitle


\begin{frame}{Introduction}
\begin{itemize}
    \item Analysts' earnings forecasts provide valuable insights into the expectations and uncertainty surrounding firms' financial performance.
    \item Forecast dispersion, measured as the disagreement among analysts, can serve as a proxy for uncertainty about a firm's future earnings.
    \item This study explores the properties of forecast dispersion and forecast errors using a comprehensive dataset of Japanese firms from 1985 to 2023.
    \item We investigate the determinants of forecast dispersion and errors, as well as their relationships with macroeconomic and financial market indicators.
\end{itemize}
\end{frame}

\begin{frame}{Information and Macroeconomic Fluctuations}
\begin{itemize}
    \item Information plays a crucial role in shaping macroeconomic fluctuations and the business cycle.
    \item The dispersion of beliefs and expectations among economic agents can influence their decision-making and aggregate outcomes.
    \item Analysts' earnings forecasts provide a unique lens to study the role of information and uncertainty in the economy.
    \item By examining the properties of forecast dispersion and errors across firms and over time, we can gain insights into the relationship between firm-level uncertainty and macroeconomic fluctuations.
\end{itemize}
\end{frame}

\begin{frame}{Overview}
\begin{itemize}
    \item We construct a comprehensive dataset by combining earnings forecasts from IBES, stock prices from NikkeiNeeds, and financial information from the Development Bank of Japan.
    \item Our dataset covers a wide range of Japanese firms from 1985 to 2023, allowing us to study the evolution of forecast dispersion and errors over time.
    \item We define key variables such as forecast dispersion, forecast errors, and the number of analysts covering each firm.
    \item Our analysis reveals several interesting patterns:
    \begin{itemize}
        \item Forecast dispersion and errors are positively correlated, indicating that greater disagreement among analysts is associated with larger forecast errors.
        \item Forecast dispersion tends to be smaller for larger firms with more analyst coverage.
        \item The number of analysts covering a firm is positively related to its size, age, and stock volatility.
    \end{itemize}
    \item We also construct time-series indices of forecast dispersion and errors and explore their relationships with macroeconomic and financial market indicators.
    \item Our findings suggest that forecast dispersion and errors are countercyclical and negatively correlated with stock market performance.
\end{itemize}
\end{frame}


\begin{frame}{IBES (Institutional Brokers' Estimate System) Data }
  \begin{itemize}
    \item IBES Data is a database that collects and compiles financial analyst earnings estimates and recommendations
    \item Provided by Refinitiv (available on WRDS)
    \item Provides a comprehensive view of analyst expectations and consensus estimates:
    \begin{itemize}
      \item Earnings per share (EPS) estimates
      \item Buy/hold/sell recommendations
      \item Price targets
    \end{itemize}
    \item Covers a wide range of companies across various industries and regions
    \item Data is available at different frequencies (e.g., monthly, quarterly, annual)
  \end{itemize}
\end{frame}

\begin{frame}{Dataset construction}

\begin{enumerate}
    \item IBES: EPS forecasts by analysts
    \item NikkeiNeeds: Index and individual stock prices
    \item DBJ (Development Bank of Japan): Financial information
\end{enumerate}

\vspace{1em}
Dataset A
\begin{itemize}
    \item IBES and NikkeiNeeds
    \item consensus forecasts, forecast dispersion, forecast errors, stock prices and volatility
    \item Monthly (1985 January --- 2023 June)
\end{itemize}

\vspace{0.5em}
Dataset B
\begin{itemize}
    \item IBES, NikkeiNeeds, and DBJ
    \item Dataset A plus financial information of individual firms
    \item Yearly  (1985 --- 2023)
\end{itemize}
\end{frame}

\begin{frame}{IBES being matched with NikkeiNeeds}
    \centering
    \includegraphics[width=1.0\linewidth]{graph/merge_nikkei_counts.png}
\end{frame}



\begin{frame}{Variable construction 1}

EPS (earnings per share) forecast: $f_{hijt}$
\vspace{1em}
\begin{itemize}
    \item forecast horizon $h$ (e.g. six months ahead of Fiscal Year End)
    \item forecast by analyst $i$
    \item forecast about firm $j$
    \item forecast time $t$ (e.g. 2005 January) 
\end{itemize}

\vspace{1em}

\end{frame}


\begin{frame}{Variable construction 2}

Forecast dispersion:

\begin{equation}
Fdis_{hjt} = \frac{\sigma_{hjt}}{| \bar{f}_{hjt} |} = \frac{\sqrt{\frac{1}{N_{hjt}-1} \sum_{i=1}^{N_{hjt}} (f_{hijt} - \bar{f}_{hjt})^2}}{| \bar{f}_{hjt} |}
\end{equation}

\vspace{2em}
--- $N_{hjt}$ is the number of analysts who made forecasts for the given forecast horizon $h$, firm $j$, and time $t$.

\vspace{1em}
--- $\bar{f}_{hjt}$ is the median of the forecasts across analysts for the given $h$, $j$, and $t$.

\end{frame}


\begin{frame}{Variable construction 3}

Forecast errors:

\begin{equation}
FE_{hjt} = | log(\frac{e_{jy(t)}}{\bar{f}_{hjt}}) |
\end{equation}

\vspace{1em}
--- $N_{hjt}$ is the number of analysts who made forecasts for the given forecast horizon $h$, firm $j$, and time $t$.

\vspace{1em}
--- $e_{jy(t)}$ is the realised earnings per share for the firm $j$ and year $y(t)$ of time $t$. For instance, $y(t)$ for $t=(2004 \text{~January})$ would be $2004$.


\end{frame}



\begin{frame}{Toyota}
    
\end{frame}

\begin{frame}{Uniqlo}
    
\end{frame}

\begin{frame}{Rakuten}
    
\end{frame}

\begin{frame}{Descriptive Statistics: Panel A}

\begin{table}
\centering
\input{table/desc_stats_a.tex}
\label{tab:desc_stats}
\end{table}

\end{frame}

\begin{frame}{Descriptive Statistics: Panel B}
\small
\begin{table}
\centering
\input{table/desc_stats_b.tex}
\label{tab:desc_stats}
\end{table}

\begin{itemize}
\item The mean sales is 420 billion Yen. The mean market cap. is 270 billion yen.
\end{itemize}
\end{frame}


\begin{frame}{Forecasts start to be released 10 month ahead}

      \begin{figure}
        \centering
        \includegraphics[width=\textwidth]{graph/numfirm_h.png}
        \caption{Forecasts start to be released 10 month ahead}
        \label{fig:figure1}
      \end{figure}
      
\end{frame}

\begin{frame}{Forecasters get better informed gradually}

  \begin{columns}
    \begin{column}{0.5\textwidth}
      \begin{figure}
        \centering
        \includegraphics[width=\textwidth]{graph/meanfdis_h.png}
        \caption{Forecast Dispersion}
        \label{fig:figure1}
      \end{figure}
    \end{column}
    
    \begin{column}{0.5\textwidth}
      \begin{figure}
        \centering
        \includegraphics[width=\textwidth]{graph/meanFElog_h.png}
        \caption{Forecast Errors}
        \label{fig:figure2}
      \end{figure}
    \end{column}
  \end{columns}
  
\end{frame}

\begin{comment}
    
\begin{frame}
	\frametitle{Forecast Dispersion (winsorized at 1\%)}
	\centering
	\includegraphics[width=1.0\linewidth]{graph/meanfdiscv_m_winsor.png}

\end{frame}
\end{comment}


\begin{frame}{Number of forecasts matter?}

  \begin{columns}
    \begin{column}{0.5\textwidth}
      \begin{figure}
        \centering
        \includegraphics[width=\textwidth]{graph/StdevNumest.png}
        \caption{Standard Deviation of Forecasts}
        \label{fig:figure1}
      \end{figure}
    \end{column}
    
    \begin{column}{0.5\textwidth}
      \begin{figure}
        \centering
        \includegraphics[width=\textwidth]{graph/ActualNmest.png}
        \caption{Realised EPS}
        \label{fig:figure2}
      \end{figure}
    \end{column}
  \end{columns}
  
\end{frame}

\begin{frame}{Information Spillover }

  \begin{columns}
    \begin{column}{0.5\textwidth}
      \begin{figure}
        \centering
        \includegraphics[width=\textwidth]{graph/FdisCVnumest.png}
        \caption{Forecast Dispersion}
        \label{fig:figure1}
      \end{figure}
    \end{column}
    
    \begin{column}{0.5\textwidth}
      \begin{figure}
        \centering
        \includegraphics[width=\textwidth]{graph/FElogNumest.png}
        \caption{Forecast Errors}
        \label{fig:figure2}
      \end{figure}
    \end{column}
  \end{columns}
  
\end{frame}


\begin{frame}{Forecast dispersion }

\begin{table}[!htbp]
\centering
\resizebox{\linewidth}{!}{ % begin threeparttable here
\begin{threeparttable} % tabular file input, then table notes
\input{table/reg_T01.tex}
\begin{tablenotes}
%\scriptsize 
\footnotesize
\item 
\end{tablenotes}
\end{threeparttable}
} % end threeparttable here
\end{table}

\end{frame}


\begin{frame}{Forecast errors }

\begin{table}[!htbp]
\centering
\resizebox{\linewidth}{!}{ % begin threeparttable here
\begin{threeparttable} % tabular file input, then table notes
\input{table/reg_T02.tex}
\begin{tablenotes}
%\scriptsize 
\footnotesize
\item 
\end{tablenotes}
\end{threeparttable}
} % end threeparttable here
\end{table}

\end{frame}


\begin{frame}{Number of estimators }

\begin{table}[!htbp]
\centering
\resizebox{\linewidth}{!}{ % begin threeparttable here
\begin{threeparttable} % tabular file input, then table notes
\input{table/reg_T03.tex}
\begin{tablenotes}
%\scriptsize 
\footnotesize
\item 
\end{tablenotes}
\end{threeparttable}
} % end threeparttable here
\end{table}

\end{frame}


\begin{frame}{Key Findings so far}
\begin{enumerate}

    \item Forecast dispersion and forecast errors covary.
    \item Major forecasts start 10 month ahead of the end of year.
    \item Smaller forecast dispersion for larger firms with more analysts covering.
    \item Large, old, and more volatile stocks have more coverage analysts. 

\end{enumerate}
\end{frame}

\begin{frame}{Time Series Index}

\begin{itemize}
    \item Forecast dispersion: $Fdis_{hjt}$
    \item Forecast errors: $FE_{hjt}$
\end{itemize}

\vspace{1em}
We take the annual mean for each firm $j$ in each year $t$. We then take the cross sectional mean for each year $t$. 


\end{frame}

\begin{frame}
	\frametitle{Forecast Dispersion (full sample)}
	\centering
	\includegraphics[width=1.0\linewidth]{graph/meanfdiscv_m.png}

\end{frame}



\begin{frame}
	\frametitle{Economic Policy Uncertainty in Japan}
	\centering
	\includegraphics[width=1.0\linewidth]{graph/mean_Fdis_CV_JPNEPUINDXM.png}

\end{frame}

\begin{comment}
\begin{frame}
	\frametitle{Industrial Production (total manufacturing)}
	\centering
	\includegraphics[width=1.0\linewidth]{graph/mean_Fdis_CV_JPNPRMNTO01GYSAM.png} 

\end{frame}
\end{comment}


\begin{frame}
	\frametitle{Nikkei 225 (monthly average)}
	\centering
	\includegraphics[width=1.0\linewidth]{graph/mean_Fdis_CV_NIKKEI225.png}

\end{frame}

\begin{frame}
	\frametitle{Moving average Fdis}
	\centering
	\includegraphics[width=1.0\linewidth]{graph/mean_Fdis_CV_MA.png}
\end{frame}

\begin{frame}
	\frametitle{Moving average FE}
	\centering
	\includegraphics[width=1.0\linewidth]{graph/mean_FE_pct_MA.png}
\end{frame}


\begin{frame}[fragile]
	\frametitle{Cross Correlation}
	\begin{table}[!htbp]
	\centering
	\resizebox{\linewidth}{!}{
	\begin{threeparttable}
	\input{table/cross_correlation.tex}
	\begin{tablenotes}
	\footnotesize
	\item
	\end{tablenotes}
	\end{threeparttable}
	}
	\end{table}
\end{frame}

\begin{frame}[fragile]
\frametitle{Uncertainty rises when stock price falls}

\footnotesize
\begin{table}[!htbp]
\centering
\resizebox{\linewidth}{!}{
\begin{threeparttable}
\input{table/reg_ts_1.tex}
\begin{tablenotes}
\footnotesize
\item
\end{tablenotes}
\end{threeparttable}
}
\end{table}
\begin{table}[!htbp]
\centering
\resizebox{\linewidth}{!}{
\begin{threeparttable}
\input{table/reg_ts_2.tex}
\begin{tablenotes}
\footnotesize
\item
\end{tablenotes}
\end{threeparttable}
}
\end{table}
\end{frame}


\begin{frame}{Summary}

\begin{itemize}
    \item Time series index based on forecast dispersion and forecast errors are correlated with other popular uncertainty measures like EPU (economic policy uncertainty index).
    \item Each spike appears to be related to an event that created uncertainty.
    \item Moving average is also informative to gauge the level of uncertainty in the medium run. 
    \item They are all countercyclical. For instance, they are negatively correlated with industrial production index and Nikkei stock index.
    
\end{itemize}
    
\end{frame}
    
\end{frame}











\end{document}